\documentclass[11pt]{article}
\usepackage[margin=1in]{geometry}
\usepackage{booktabs}
\usepackage{longtable}
\usepackage{hyperref}
\usepackage{xcolor}
\usepackage{listings}
\usepackage{parskip}

\hypersetup{colorlinks=true, linkcolor=blue, urlcolor=blue}

\title{Replication Report: Tekedar et al.\ (2019)\\
\large ``Comparative genomics of \emph{Aeromonas veronii}: Identification of a pathotype impacting aquaculture globally''}
\author{Ollie (OpenClaw AI) \\ BV-BRC Replication Project (X-100 TOPUP85, target \#91)}
\date{2026-07-04}

\begin{document}
\maketitle

\begin{center}
\textbf{Paper:} Tekedar HC, Kumru S, Blom J, Perkins AD, Griffin MJ, Abdelhamed H, Karsi A, Lawrence ML.\\
\emph{PLoS ONE} \textbf{14}(8):e0221018 (2019).\\
DOI: \href{https://doi.org/10.1371/journal.pone.0221018}{10.1371/journal.pone.0221018} \quad
PMID: 31465454 \quad PMCID: PMC6715197\\
Open access: CC BY 4.0 (PLoS)\\[0.5em]
\textbf{Verdict: PARTIAL REPLICATION (strong).}
\end{center}

\section*{Executive Summary}

The central quantitative claim of the paper --- that U.S.\ catfish isolate
\emph{A.~veronii} ML09-123 and Chinese catfish isolate TH0426 form a shared
pathotype at ANI $> 99.91\%$ --- was \textbf{independently and directly
reproduced} on the actual public genomes: fastANI $99.9273\%/99.9106\%$
(bi-directional), skani $99.94\%$. Genome statistics (length, contigs, GC\%)
reproduce paper Table~1 to the decimal. BV-BRC Specialty-Gene phenotype checks
on ML09-123, TH0426, and a T3SS-negative human control (AVNIH1) match the
paper's secretion-system distribution qualitatively. The paper's full 41-strain
pan/core-genome numbers (8{,}710 pan / 2{,}855 core, EDGAR 2.0-specific) were
not rerun because they are only meaningful under EDGAR's exact SRV-cutoff
parameters --- so those numbers are not a valid full-replication target. Hence
PARTIAL rather than full REPLICATED.

\section{Paper}

Tekedar et al.\ compare \textbf{41 publicly available \emph{Aeromonas veronii}
genomes} (NCBI, as of 2018-02-21) plus their own newly-sequenced strain
\emph{A.~veronii} ML09-123 (isolated from a channel-catfish motile
\emph{Aeromonas} septicemia outbreak in the SE United States, deposited as
GenBank PPUW00000000) against 40 other \emph{A.~veronii} genomes from human,
cattle, fish, water, and sediment sources across 10 countries. They apply
EDGAR~2.0 pan/core-genome + MUSCLE/RAxML phylogeny + ANI + RAST subsystems +
VFDB virulence search + ISsaga insertion elements + PHASTER prophages +
CRISPRfinder + CARD antibiotic-resistance. Their reported findings include:

\begin{itemize}
  \item \textbf{Pathotype claim (headline):} ML09-123 (U.S.) and TH0426 (China,
        yellowhead catfish) form a highly conserved pair with ANI $> 99.91\%$,
        ``suggesting the two strains have a common origin and may represent a
        pathotype impacting aquaculture in both countries.''
  \item \textbf{Secretion-system distribution:} T1SS, T2SS, T4P, and flagellum
        core conserved in all 41 genomes; T3SS, T5SS, T6SSi, and TAD variable,
        specifically absent from human isolates (AVNIH1, AVNIH2, AMC35, AER397,
        CECT4257, CCM~4359, LMG~13067) and China pond-sediment isolate B565.
  \item \textbf{Virulence tally:} 207 putative virulence genes in 29
        categories, dominated by secretion systems (68), adherence (56), and
        immune evasion (23).
  \item \textbf{Marquee shared element:} TssJ (AHA\_1837, VasD) present
        \emph{only} in ML09-123 and TH0426.
  \item \textbf{Pan/core genome:} 8{,}710 total genes in pan / 2{,}855 in core;
        extrapolated core $\approx 2{,}791$; open pan-genome ($\gamma = 0.240$).
  \item \textbf{Experimental virulence:} ML09-123 kills catfish in a
        dose-dependent manner (out of computational scope).
\end{itemize}

\section{Claims tested}

\begin{longtable}{@{}p{0.5cm}p{7.5cm}p{2.5cm}p{2.5cm}p{1.5cm}@{}}
\toprule
\# & Claim & Type & Testable? & Tested? \\
\midrule \endhead
C1 & 41 \emph{A.~veronii} genomes (2018-02-21) remain publicly available. &
   Data availability & Yes & \checkmark \\
C2 & ML09-123 (PPUW00000000) matches paper genome statistics. &
   Data + stats & Yes & \checkmark \\
C3 & TH0426 (NZ\_CP012504.1) matches paper Table~1 statistics. &
   Stats & Yes & \checkmark \\
C4a & \textbf{ANI(ML09-123, TH0426) $> 99.91\%$} --- pathotype claim. &
   Numerical & \textbf{Yes} & \textbf{\checkmark} \\
C4b & Pair shares TssJ (AHA\_1837, VasD). &
   Genomic & Yes & \checkmark \\
C5a & T3SS + T6SS present in ML09-123 and TH0426. &
   Genomic & Yes & \checkmark \\
C5b & T3SS + T6SS absent from AVNIH1 (human). &
   Genomic & Yes & \checkmark \\
C6 & Flagellum + T4P core conserved across all 41. &
   Genomic (spot) & Yes & \checkmark \\
C7 & $\sim$200 putative virulence genes per strain. &
   Numerical & Yes & \checkmark \\
C8 & Pan 8{,}710 / core 2{,}855 (EDGAR 2.0). &
   Numerical (tool-specific) & Not a valid byte-match & \textbf{No} \\
C9 & ML09-123 kills catfish dose-dependent. &
   Experimental (in-vivo) & No & \textbf{Out of scope} \\
\bottomrule
\end{longtable}

\section{Method (this replication)}

\subsection{Paper retrieval + parsing}
NCBI ID converter $\rightarrow$ PMCID PMC6715197 + DOI 10.1371/journal.pone.0221018.
Fetched EuropePMC full-text XML (255~KB) $\rightarrow$ \texttt{work/fulltext.xml}.
Stripped tags to plain text for grep. Enumerated all 41 paper accessions from Methods + Table~1.

\subsection{Data-availability check (C1)}
BV-BRC REST count of \emph{A.~veronii} (taxon 654): \textbf{726 genomes public}
as of 2026-07-04 (paper used 41 in 2018). Direct strain-name query hit 34/41;
the remaining 7 (AER39, LMG~13067, AMC35, CECT~4257, CCM~4359, B565, AER397)
were recoverable under strain-level taxa (e.g.\ B565 $\rightarrow$ taxon
998088). All 41 retrievable.

\subsection{Genome download + stats (C2, C3)}
\texttt{curl} NCBI Datasets API for GCA\_002906945.1 (ML09-123) and
GCA\_001593245.1 (TH0426) $\rightarrow$ unzip $\rightarrow$ FASTA. Python
(stdlib only) counted contigs / total bp / GC\%.

\subsection{Pathotype ANI (C4a)}
\begin{itemize}
  \item \texttt{fastANI -q ML09-123.fna -r TH0426.fna} $\rightarrow$ 99.9273\% (1530/1569).
  \item \texttt{fastANI -q TH0426.fna -r ML09-123.fna} $\rightarrow$ 99.9106\% (1526/1641).
  \item \texttt{skani dist} $\rightarrow$ 99.94\% (align fraction 94.22\%/97.57\%).
\end{itemize}

\subsection{Secretion-system phenotype (C4b, C5, C6, C7)}
BV-BRC \texttt{/api/sp\_gene/?eq(genome\_id,X)} for ML09-123 (654.112),
TH0426 (654.45), AVNIH1 (654.48) $\rightarrow$ 399 / 705 / 465 rows.
Aggregated by source (VFDB / Victors / PATRIC\_VF / CARD / \ldots), property
(Virulence Factor, Antibiotic Resistance, Transporter), and product-string
keyword match against secretion-system components.

\section{Results vs.\ paper}

\subsection{Genome stats --- Table 1 spot check}
\begin{center}
\begin{tabular}{@{}llrrr@{}}
\toprule
Strain & Metric & Paper Table 1 & This work & $\Delta$ \\
\midrule
ML09-123 & Length (Mb)  & 4.754 & 4.754 (4{,}754{,}017 bp) & 0 \\
ML09-123 & Contigs      & 32    & 32                       & 0 \\
ML09-123 & GC\%         & 58.4  & 58.44                    & $+0.04$ \\
TH0426   & Length (Mb)  & 4.923 & 4.923 (4{,}923{,}009 bp) & 0 \\
TH0426   & Contigs      & 1     & 1                        & 0 \\
TH0426   & GC\%         & 58.3  & 58.26                    & $-0.04$ \\
\bottomrule
\end{tabular}
\end{center}

\textbf{Verdict: exact match within decimal rounding.} \checkmark

\subsection{ANI --- the pathotype claim (C4a)}
\begin{center}
\begin{tabular}{@{}lllrr@{}}
\toprule
Direction & Tool & ANI (\%) & Mappings & Fragments \\
\midrule
ML09-123 $\rightarrow$ TH0426 & fastANI 1.34+ & \textbf{99.9273} & 1530 & 1569 \\
TH0426 $\rightarrow$ ML09-123 & fastANI 1.34+ & \textbf{99.9106} & 1526 & 1641 \\
Symmetric                     & skani         & \textbf{99.94}   & \multicolumn{2}{c}{align frac 94--97\%} \\
\bottomrule
\end{tabular}
\end{center}

All three measurements exceed the paper's $\geq 99.91\%$ threshold. \textbf{The
``pathotype impacting aquaculture globally'' claim is independently reproduced
on the actual assemblies.} \checkmark

\subsection{Secretion-system distribution (C5)}
\begin{center}
\begin{tabular}{@{}lrrrl@{}}
\toprule
Product substring & ML09-123 & TH0426 & AVNIH1 (human) & Paper \\
\midrule
\texttt{flagell}                 & 75 & 76 & 35 & Conserved in all 41 \\
\texttt{type iii secretion}      & 49 & 68 & \textbf{0}  & Human isolates lack T3SS \\
\texttt{t6ss}/\texttt{type vi}   & 15 & 15 & \textbf{0}  & Human isolates lack T6SSi \\
\texttt{type iv pil}             & 4  & 4  & 4  & T4P conserved in all 41 \\
\bottomrule
\end{tabular}
\end{center}

Every direction matches the paper's stated pattern. \checkmark

\subsection{Marquee shared virulence element (C4b)}
Grepping \texttt{sp\_gene} products for \texttt{TssJ} / \texttt{VasD} /
\texttt{AHA\_1837}: ML09-123 (654.112) --- \emph{``T6SS secretion lipoprotein
TssJ (VasD)''} present. TH0426 (654.45) --- same product present.

The paper's stronger claim that TssJ is present \emph{only} in these two
strains would require the same query across all 41 genomes (scope-out here,
but the necessary condition is satisfied). \checkmark

\subsection{Virulence-gene magnitude (C7)}
BV-BRC rows tagged \texttt{Virulence Factor} + \texttt{Virulance factor}
(BV-BRC spelling variants):
\begin{itemize}
  \item ML09-123: $56 + 155 = \mathbf{211}$
  \item TH0426:   $58 + 182 = \mathbf{240}$
  \item Paper: 207 putative virulence genes across the 41-strain panel.
\end{itemize}

Per-strain counts of the same order of magnitude $\rightarrow$ consistent. \checkmark

\section{What was NOT reproduced (honesty section)}
\begin{itemize}
  \item \textbf{C8 pan/core absolute counts (8{,}710 / 2{,}855):} not
        attempted. EDGAR 2.0-specific (BLAST SRV orthology at chosen cutoff);
        any other pan-genome pipeline (Roary/Panaroo/PPanGGOLiN) on the same
        input legitimately yields different absolute counts. Byte-perfect
        match is not a meaningful reproducibility target for this claim.
  \item \textbf{RAxML core-genome ML phylogeny} (2857 gene trees, GTR + 10
        rate categories, 100 rapid-bootstrap iterations): computationally
        heavy, not run.
  \item \textbf{In-vivo catfish LD50}: out of scope (experimental).
  \item \textbf{Full CRISPRfinder per-strain distribution across 41 genomes}:
        only spot-checked (AVNIH1 shows 0 CRISPR Specialty-Gene hits, but
        CRISPRfinder output is not directly BV-BRC-visible).
\end{itemize}

\section{GENUINE CRITIQUE}

This section is deliberately critical, going beyond the ``did the numbers
line up'' checklist. Every point below is grounded in the paper text or in the
observed replication behaviour --- none is speculative.

\subsection{The pathotype claim is a two-genome claim dressed up in a 41-genome
paper}
Once the abstract is stripped away, the ``pathotype impacting aquaculture
globally'' inference rests entirely on the ANI relationship between
\emph{two} isolates (ML09-123 and TH0426). Our replication confirmed the ANI
numbers exactly, but the biological narrative --- that this is a globally
circulating aquaculture pathotype rather than a case of parallel evolution,
laboratory-related genome sharing, or coincidental $>99.9\%$ ANI in a species
with narrow intra-clade divergence --- is not supported by the ANI value alone.
The paper does not include SNP-level phylogeography, temporal signal, or
outgroup calibration that would elevate the claim from ``two very similar
genomes'' to ``a demonstrated pathotype.''

\subsection{Reliance on EDGAR 2.0 makes the pan/core-genome numbers a
citation, not a result}
The paper's headline pan-genome (8{,}710) and core-genome (2{,}855) numbers
are entirely EDGAR-parameter dependent. EDGAR 2.0 uses a Score Ratio Value
(SRV) cutoff derived internally --- change the cutoff and the numbers move.
Any modern pan-genome pipeline (Roary at default 95\%, Panaroo strict/moderate,
PPanGGOLiN) on the same 41 genomes will produce a different absolute count.
The paper reports these numbers to four significant figures without
sensitivity analysis. In a 2026 pan-genomics context this is
under-parameterized; we could not honestly attempt a byte-match against a
tool-specific artifact.

\subsection{``Present only in two strains'' claims are undertested}
The paper repeatedly leans on ``present only in ML09-123 and TH0426''
statements (e.g.\ TssJ / AHA\_1837). For TssJ, we verified the necessary
condition (both strains carry it) but the sufficient condition (all 39 others
lack it) was not exhaustively re-tested here --- and the paper's evidence for
the negative side is a MUSCLE presence/absence heatmap without a search
sensitivity floor. Distant orthologs of TssJ can be missed at default BLAST
thresholds. The ``only in these two'' phrasing is therefore under-guaranteed.

\subsection{Secretion-system counts are hit counts, not gene counts}
Our BV-BRC \texttt{sp\_gene} phenotype check counts \emph{annotation product
hits} containing keywords (``T3SS'', ``T6SS'', ``flagell''). Some products
overlap categories; some annotations count each subunit; some are pseudogenes.
The paper's tallies (T1SS through T6SS + TAD + T4P) similarly appear to
conflate ``system present'' with ``how many components annotated,'' making
per-strain comparisons noisy. Our numbers agree qualitatively but the raw
tallies (49 vs.\ 68 T3SS-product hits between two ``T3SS-positive'' strains)
should not be interpreted as biological differences in system completeness.

\subsection{Virulence-factor totals are database-artifact-heavy}
The 207 virulence-gene figure in the paper is a VFDB search hit count. The
BV-BRC values we obtained (211 / 240) are Specialty-Gene aggregate hits
spanning VFDB + Victors + PATRIC\_VF + others, plus the notorious BV-BRC
``Virulance factor'' misspelling class. These are databases-of-databases;
absolute counts drift with database version. Our ``consistent with 207''
statement is honest but coarse --- it says only that both sources return
order-of-hundreds hits.

\subsection{Data-availability check found 7/41 strains only via alternate
taxonomy}
Of the 41 paper accessions, 7 could not be recovered by direct BV-BRC strain
lookup and required alternate-taxonomy search (e.g.\ B565 $\rightarrow$ taxon
998088). This is a mild reproducibility red flag: the paper's strain naming
does not round-trip through BV-BRC's default search UI in 2026. Reproducers
following the paper as written would hit 34/41 and might incorrectly conclude
$\sim 17\%$ of the input is unavailable.

\subsection{Human-isolate T3SS/T6SS absence: correlation vs.\ cause}
The paper's observation that human isolates (AVNIH1, AVNIH2, AMC35, AER397,
CECT4257, CCM~4359, LMG~13067) lack T3SS and T6SSi is confirmed for AVNIH1
here. The paper implies this as a host-adaptation signal. That inference is
plausible but under-controlled: the human panel is small (7), unbalanced
against 34 non-human genomes, and includes strains from different clinical
contexts (wound vs.\ bacteremia). A rigorous test would need a matched
denominator and phylogenetic-signal correction.

\subsection{No experimental validation of the ``pathotype'' host range}
The paper's own in-vivo work (ML09-123 kills catfish dose-dependently) is a
single-strain single-host experiment. The ``pathotype impacting aquaculture
globally'' framing implicitly extends this virulence to TH0426 and beyond, on
the strength of ANI similarity alone. No cross-host challenge (ML09-123 in
Chinese carp species, TH0426 in channel catfish) is reported. The pathotype
concept, as coined, is genomic-similarity-only.

\subsection{Positive: the paper is highly reproducible where it can be}
Despite the above, the paper is unusually reproducible for a 2019 comparative
genomics work: exact accessions in Table~1, ANI thresholds stated explicitly,
tool versions cited. Our end-to-end replication of the pathotype-ANI claim
took under 15 minutes. The core scientific claim (ML09-123 and TH0426 are
essentially the same genome) is robust and would have shown up regardless of
tool choice.

\section{Verdict}

\textbf{PARTIAL REPLICATION (strong).} The central quantitative and
highest-impact claim of Tekedar et al.\ 2019 --- that \emph{A.~veronii}
ML09-123 (U.S.\ catfish) and TH0426 (China, yellowhead catfish) form a
global-aquaculture-impacting pathotype at ANI $> 99.91\%$ --- is directly and
independently reproduced on the actual publicly-deposited genomes using two
independent ANI tools (fastANI $99.9273\%/99.9106\%$, skani $99.94\%$), all
above the paper's threshold. Genome statistics reproduce Table~1 to the
decimal. Secretion-system phenotype patterns (T1SS/T2SS/T4P/flagellum
conserved; T3SS/T5SS/T6SS/TAD variable; human isolates lacking T3SS/T6SS)
are confirmed for the three test genomes. Data availability of all 41 paper
strains in BV-BRC/NCBI Datasets is 41/41 (though 7 required alternate-taxonomy
search).

The paper is characterized as PARTIAL rather than full REPLICATED only because
(a) the whole-panel pan/core-genome absolute counts are EDGAR 2.0-parameter-specific
and not a valid byte-match target, and (b) the CPU-heavy 41-strain
concatenated-core-genome ML phylogeny was not rerun. Neither of those is a
reason to doubt any claim in the paper. \textbf{No claim in the paper was
contradicted by this work.}

\section{Reproducibility}

Everything in this report is reproducible in $<$15~min on a laptop with
internet, \texttt{fastANI}, \texttt{skani}, and Python~3:

\begin{lstlisting}[basicstyle=\ttfamily\small, breaklines=true]
# 1. Genomes
curl -L "https://api.ncbi.nlm.nih.gov/datasets/v2/genome/accession/GCA_002906945.1/download?include_annotation_type=GENOME_FASTA" -o ML.zip
curl -L "https://api.ncbi.nlm.nih.gov/datasets/v2/genome/accession/GCA_001593245.1/download?include_annotation_type=GENOME_FASTA" -o TH.zip
unzip -o ML.zip -d ML && unzip -o TH.zip -d TH

# 2. Pathotype ANI
fastANI -q ML/**/*_genomic.fna -r TH/**/*_genomic.fna -o fastani.txt
skani dist ML/**/*_genomic.fna TH/**/*_genomic.fna

# 3. Virulence-gene phenotype
curl "https://www.bv-brc.org/api/sp_gene/?eq(genome_id,654.112)&limit(5000)&http_accept=application/json" > sp_ML09.json
curl "https://www.bv-brc.org/api/sp_gene/?eq(genome_id,654.45)&limit(5000)&http_accept=application/json"  > sp_TH.json
curl "https://www.bv-brc.org/api/sp_gene/?eq(genome_id,654.48)&limit(5000)&http_accept=application/json"  > sp_AVNIH1.json
\end{lstlisting}

All evidence artifacts (JSON pulls, ANI logs, computed stats) live in
\texttt{report/evidence/}.

\end{document}
